% Options for packages loaded elsewhere
\PassOptionsToPackage{unicode}{hyperref}
\PassOptionsToPackage{hyphens}{url}
\documentclass[
]{article}
\usepackage{xcolor}
\usepackage[margin=1in]{geometry}
\usepackage{amsmath,amssymb}
\setcounter{secnumdepth}{-\maxdimen} % remove section numbering
\usepackage{iftex}
\ifPDFTeX
  \usepackage[T1]{fontenc}
  \usepackage[utf8]{inputenc}
  \usepackage{textcomp} % provide euro and other symbols
\else % if luatex or xetex
  \usepackage{unicode-math} % this also loads fontspec
  \defaultfontfeatures{Scale=MatchLowercase}
  \defaultfontfeatures[\rmfamily]{Ligatures=TeX,Scale=1}
\fi
\usepackage{lmodern}
\ifPDFTeX\else
  % xetex/luatex font selection
\fi
% Use upquote if available, for straight quotes in verbatim environments
\IfFileExists{upquote.sty}{\usepackage{upquote}}{}
\IfFileExists{microtype.sty}{% use microtype if available
  \usepackage[]{microtype}
  \UseMicrotypeSet[protrusion]{basicmath} % disable protrusion for tt fonts
}{}
\makeatletter
\@ifundefined{KOMAClassName}{% if non-KOMA class
  \IfFileExists{parskip.sty}{%
    \usepackage{parskip}
  }{% else
    \setlength{\parindent}{0pt}
    \setlength{\parskip}{6pt plus 2pt minus 1pt}}
}{% if KOMA class
  \KOMAoptions{parskip=half}}
\makeatother
\usepackage{color}
\usepackage{fancyvrb}
\newcommand{\VerbBar}{|}
\newcommand{\VERB}{\Verb[commandchars=\\\{\}]}
\DefineVerbatimEnvironment{Highlighting}{Verbatim}{commandchars=\\\{\}}
% Add ',fontsize=\small' for more characters per line
\usepackage{framed}
\definecolor{shadecolor}{RGB}{248,248,248}
\newenvironment{Shaded}{\begin{snugshade}}{\end{snugshade}}
\newcommand{\AlertTok}[1]{\textcolor[rgb]{0.94,0.16,0.16}{#1}}
\newcommand{\AnnotationTok}[1]{\textcolor[rgb]{0.56,0.35,0.01}{\textbf{\textit{#1}}}}
\newcommand{\AttributeTok}[1]{\textcolor[rgb]{0.13,0.29,0.53}{#1}}
\newcommand{\BaseNTok}[1]{\textcolor[rgb]{0.00,0.00,0.81}{#1}}
\newcommand{\BuiltInTok}[1]{#1}
\newcommand{\CharTok}[1]{\textcolor[rgb]{0.31,0.60,0.02}{#1}}
\newcommand{\CommentTok}[1]{\textcolor[rgb]{0.56,0.35,0.01}{\textit{#1}}}
\newcommand{\CommentVarTok}[1]{\textcolor[rgb]{0.56,0.35,0.01}{\textbf{\textit{#1}}}}
\newcommand{\ConstantTok}[1]{\textcolor[rgb]{0.56,0.35,0.01}{#1}}
\newcommand{\ControlFlowTok}[1]{\textcolor[rgb]{0.13,0.29,0.53}{\textbf{#1}}}
\newcommand{\DataTypeTok}[1]{\textcolor[rgb]{0.13,0.29,0.53}{#1}}
\newcommand{\DecValTok}[1]{\textcolor[rgb]{0.00,0.00,0.81}{#1}}
\newcommand{\DocumentationTok}[1]{\textcolor[rgb]{0.56,0.35,0.01}{\textbf{\textit{#1}}}}
\newcommand{\ErrorTok}[1]{\textcolor[rgb]{0.64,0.00,0.00}{\textbf{#1}}}
\newcommand{\ExtensionTok}[1]{#1}
\newcommand{\FloatTok}[1]{\textcolor[rgb]{0.00,0.00,0.81}{#1}}
\newcommand{\FunctionTok}[1]{\textcolor[rgb]{0.13,0.29,0.53}{\textbf{#1}}}
\newcommand{\ImportTok}[1]{#1}
\newcommand{\InformationTok}[1]{\textcolor[rgb]{0.56,0.35,0.01}{\textbf{\textit{#1}}}}
\newcommand{\KeywordTok}[1]{\textcolor[rgb]{0.13,0.29,0.53}{\textbf{#1}}}
\newcommand{\NormalTok}[1]{#1}
\newcommand{\OperatorTok}[1]{\textcolor[rgb]{0.81,0.36,0.00}{\textbf{#1}}}
\newcommand{\OtherTok}[1]{\textcolor[rgb]{0.56,0.35,0.01}{#1}}
\newcommand{\PreprocessorTok}[1]{\textcolor[rgb]{0.56,0.35,0.01}{\textit{#1}}}
\newcommand{\RegionMarkerTok}[1]{#1}
\newcommand{\SpecialCharTok}[1]{\textcolor[rgb]{0.81,0.36,0.00}{\textbf{#1}}}
\newcommand{\SpecialStringTok}[1]{\textcolor[rgb]{0.31,0.60,0.02}{#1}}
\newcommand{\StringTok}[1]{\textcolor[rgb]{0.31,0.60,0.02}{#1}}
\newcommand{\VariableTok}[1]{\textcolor[rgb]{0.00,0.00,0.00}{#1}}
\newcommand{\VerbatimStringTok}[1]{\textcolor[rgb]{0.31,0.60,0.02}{#1}}
\newcommand{\WarningTok}[1]{\textcolor[rgb]{0.56,0.35,0.01}{\textbf{\textit{#1}}}}
\usepackage{graphicx}
\makeatletter
\newsavebox\pandoc@box
\newcommand*\pandocbounded[1]{% scales image to fit in text height/width
  \sbox\pandoc@box{#1}%
  \Gscale@div\@tempa{\textheight}{\dimexpr\ht\pandoc@box+\dp\pandoc@box\relax}%
  \Gscale@div\@tempb{\linewidth}{\wd\pandoc@box}%
  \ifdim\@tempb\p@<\@tempa\p@\let\@tempa\@tempb\fi% select the smaller of both
  \ifdim\@tempa\p@<\p@\scalebox{\@tempa}{\usebox\pandoc@box}%
  \else\usebox{\pandoc@box}%
  \fi%
}
% Set default figure placement to htbp
\def\fps@figure{htbp}
\makeatother
\setlength{\emergencystretch}{3em} % prevent overfull lines
\providecommand{\tightlist}{%
  \setlength{\itemsep}{0pt}\setlength{\parskip}{0pt}}
\usepackage{bookmark}
\IfFileExists{xurl.sty}{\usepackage{xurl}}{} % add URL line breaks if available
\urlstyle{same}
\hypersetup{
  pdftitle={elec\_consumption},
  hidelinks,
  pdfcreator={LaTeX via pandoc}}

\title{elec\_consumption}
\author{}
\date{\vspace{-2.5em}2026-08-03}

\begin{document}
\maketitle

We have the electricity consumption in kW that we want to forecast. We
have also the outdoor temperature of the building. the data is collected
every 15min from 1/1/2010 1:15 until 2/18/2010 23:45

\begin{Shaded}
\begin{Highlighting}[]
\FunctionTok{library}\NormalTok{(readxl)}
\NormalTok{df }\OtherTok{=} \FunctionTok{read\_excel}\NormalTok{(}\StringTok{"Elec{-}train.xlsx"}\NormalTok{)}
\CommentTok{\#rename column }
\FunctionTok{names}\NormalTok{(df)[}\FunctionTok{names}\NormalTok{(df) }\SpecialCharTok{==} \StringTok{"Power (kW)"}\NormalTok{] }\OtherTok{=} \StringTok{"power"}
\FunctionTok{names}\NormalTok{(df)[}\FunctionTok{names}\NormalTok{(df) }\SpecialCharTok{==} \StringTok{"Temp (C°)"}\NormalTok{] }\OtherTok{=} \StringTok{"temp"}
\FunctionTok{names}\NormalTok{(df)[}\FunctionTok{names}\NormalTok{(df) }\SpecialCharTok{==} \StringTok{"Timestamp"}\NormalTok{] }\OtherTok{=} \StringTok{"time"}

\NormalTok{df}\SpecialCharTok{$}\NormalTok{time }\OtherTok{=} \FunctionTok{as.POSIXct}\NormalTok{(df}\SpecialCharTok{$}\NormalTok{time, }\AttributeTok{format =} \StringTok{"\%m/\%d/\%Y \%H:\%M"}\NormalTok{)}
\CommentTok{\# we work with the data from 1/1/2010 1:15 until 2/18/2010 23:45}
\NormalTok{df }\OtherTok{=}\NormalTok{ df[df}\SpecialCharTok{$}\NormalTok{time }\SpecialCharTok{\textless{}=} \FunctionTok{as.POSIXct}\NormalTok{(}\StringTok{"2/18/2010 23:45"}\NormalTok{, }\AttributeTok{format =} \StringTok{"\%m/\%d/\%Y \%H:\%M"}\NormalTok{),]}

\FunctionTok{head}\NormalTok{(df)}
\end{Highlighting}
\end{Shaded}

\begin{verbatim}
## # A tibble: 6 x 3
##   time                power  temp
##   <dttm>              <dbl> <dbl>
## 1 2010-01-01 01:15:00  165.  10.6
## 2 2010-01-01 01:30:00  152.  10.6
## 3 2010-01-01 01:45:00  147.  10.6
## 4 2010-01-01 02:00:00  154.  10.6
## 5 2010-01-01 02:15:00  154.  10.6
## 6 2010-01-01 02:30:00  159   10.6
\end{verbatim}

\begin{Shaded}
\begin{Highlighting}[]
\FunctionTok{library}\NormalTok{(xts)}
\end{Highlighting}
\end{Shaded}

\begin{verbatim}
## Loading required package: zoo
\end{verbatim}

\begin{verbatim}
## 
## Attaching package: 'zoo'
\end{verbatim}

\begin{verbatim}
## The following objects are masked from 'package:base':
## 
##     as.Date, as.Date.numeric
\end{verbatim}

\begin{Shaded}
\begin{Highlighting}[]
\NormalTok{elect }\OtherTok{=} \FunctionTok{xts}\NormalTok{(df}\SpecialCharTok{$}\NormalTok{power, }\AttributeTok{order.by =}\NormalTok{ df}\SpecialCharTok{$}\NormalTok{time)}
\FunctionTok{plot}\NormalTok{(}
\NormalTok{  elect,}
  \AttributeTok{xaxt =} \StringTok{"n"}\NormalTok{,}
  \AttributeTok{type =} \StringTok{"l"}\NormalTok{,}
  \AttributeTok{main =} \StringTok{"15{-}minute Electricity Consumption"}\NormalTok{,}
  \AttributeTok{xlab =} \StringTok{"Day"}\NormalTok{,}
  \AttributeTok{ylab =} \StringTok{"Electricity"}
\NormalTok{)}

\CommentTok{\# add daily labels}
\NormalTok{days }\OtherTok{\textless{}{-}} \FunctionTok{seq}\NormalTok{(}
  \FunctionTok{as.POSIXct}\NormalTok{(}\StringTok{"2010{-}01{-}01"}\NormalTok{),}
  \FunctionTok{as.POSIXct}\NormalTok{(}\StringTok{"2010{-}02{-}18"}\NormalTok{),}
  \AttributeTok{by =} \StringTok{"1 day"}
\NormalTok{)}

\FunctionTok{axis}\NormalTok{(}
  \DecValTok{1}\NormalTok{,}
  \AttributeTok{at =}\NormalTok{ days,}
  \AttributeTok{labels =} \FunctionTok{format}\NormalTok{(days, }\StringTok{"\%b{-}\%d"}\NormalTok{),}
  \AttributeTok{las =} \DecValTok{2}
\NormalTok{)}
\end{Highlighting}
\end{Shaded}

\pandocbounded{\includegraphics[keepaspectratio]{elec_consumption_Naro_files/elec_consumption_Naro_files/figure-latex/unnamed-chunk-2-1.pdf}}

\begin{Shaded}
\begin{Highlighting}[]
\FunctionTok{library}\NormalTok{(dplyr)}
\end{Highlighting}
\end{Shaded}

\begin{verbatim}
## 
## ######################### Warning from 'xts' package ##########################
## #                                                                             #
## # The dplyr lag() function breaks how base R's lag() function is supposed to  #
## # work, which breaks lag(my_xts). Calls to lag(my_xts) that you type or       #
## # source() into this session won't work correctly.                            #
## #                                                                             #
## # Use stats::lag() to make sure you're not using dplyr::lag(), or you can add #
## # conflictRules('dplyr', exclude = 'lag') to your .Rprofile to stop           #
## # dplyr from breaking base R's lag() function.                                #
## #                                                                             #
## # Code in packages is not affected. It's protected by R's namespace mechanism #
## # Set `options(xts.warn_dplyr_breaks_lag = FALSE)` to suppress this warning.  #
## #                                                                             #
## ###############################################################################
\end{verbatim}

\begin{verbatim}
## 
## Attaching package: 'dplyr'
\end{verbatim}

\begin{verbatim}
## The following objects are masked from 'package:xts':
## 
##     first, last
\end{verbatim}

\begin{verbatim}
## The following objects are masked from 'package:stats':
## 
##     filter, lag
\end{verbatim}

\begin{verbatim}
## The following objects are masked from 'package:base':
## 
##     intersect, setdiff, setequal, union
\end{verbatim}

\begin{Shaded}
\begin{Highlighting}[]
\FunctionTok{library}\NormalTok{(ggplot2)}


\CommentTok{\# Create date and time{-}of{-}day variables}
\NormalTok{df\_seasonal }\OtherTok{\textless{}{-}}\NormalTok{ df }\SpecialCharTok{\%\textgreater{}\%}
  \FunctionTok{mutate}\NormalTok{(}
    \AttributeTok{date =} \FunctionTok{as.Date}\NormalTok{(time),}
    \AttributeTok{time\_minutes =} \FunctionTok{as.numeric}\NormalTok{(}\FunctionTok{format}\NormalTok{(time, }\StringTok{"\%H"}\NormalTok{)) }\SpecialCharTok{*} \DecValTok{60} \SpecialCharTok{+}
                   \FunctionTok{as.numeric}\NormalTok{(}\FunctionTok{format}\NormalTok{(time, }\StringTok{"\%M"}\NormalTok{))}
\NormalTok{  )}

\CommentTok{\# Seasonal plot}
\FunctionTok{ggplot}\NormalTok{(df\_seasonal,}
       \FunctionTok{aes}\NormalTok{(}\AttributeTok{x =}\NormalTok{ time\_minutes,}
           \AttributeTok{y =}\NormalTok{ power,}
           \AttributeTok{group =}\NormalTok{ date,}
           \AttributeTok{colour =}\NormalTok{ date)) }\SpecialCharTok{+}
  \FunctionTok{geom\_line}\NormalTok{(}\AttributeTok{linewidth =} \FloatTok{0.5}\NormalTok{, }\AttributeTok{alpha =} \FloatTok{0.7}\NormalTok{) }\SpecialCharTok{+}
  \FunctionTok{scale\_x\_continuous}\NormalTok{(}
    \AttributeTok{breaks =} \FunctionTok{seq}\NormalTok{(}\DecValTok{0}\NormalTok{, }\DecValTok{1440}\NormalTok{, }\AttributeTok{by =} \DecValTok{120}\NormalTok{),}
    \AttributeTok{labels =} \FunctionTok{sprintf}\NormalTok{(}\StringTok{"\%02d:\%02d"}\NormalTok{,}
                     \FunctionTok{seq}\NormalTok{(}\DecValTok{0}\NormalTok{, }\DecValTok{1440}\NormalTok{, }\AttributeTok{by =} \DecValTok{120}\NormalTok{) }\SpecialCharTok{\%/\%} \DecValTok{60}\NormalTok{,}
                     \FunctionTok{seq}\NormalTok{(}\DecValTok{0}\NormalTok{, }\DecValTok{1440}\NormalTok{, }\AttributeTok{by =} \DecValTok{120}\NormalTok{) }\SpecialCharTok{\%\%} \DecValTok{60}\NormalTok{)}
\NormalTok{  ) }\SpecialCharTok{+}
  \FunctionTok{labs}\NormalTok{(}
    \AttributeTok{title =} \StringTok{"Daily Seasonal Pattern of Electricity Consumption"}\NormalTok{,}
    \AttributeTok{x =} \StringTok{"Time of Day"}\NormalTok{,}
    \AttributeTok{y =} \StringTok{"Electricity"}\NormalTok{,}
    \AttributeTok{colour =} \StringTok{"Date"}
\NormalTok{  ) }\SpecialCharTok{+}
  \FunctionTok{theme\_minimal}\NormalTok{() }\SpecialCharTok{+}
  \FunctionTok{theme}\NormalTok{(}
    \AttributeTok{axis.text.x =} \FunctionTok{element\_text}\NormalTok{(}\AttributeTok{angle =} \DecValTok{45}\NormalTok{, }\AttributeTok{hjust =} \DecValTok{1}\NormalTok{)}
\NormalTok{  )}
\end{Highlighting}
\end{Shaded}

\pandocbounded{\includegraphics[keepaspectratio]{elec_consumption_Naro_files/elec_consumption_Naro_files/figure-latex/unnamed-chunk-3-1.pdf}}

We can see that there is daily pattern in this dataset: - mid night
until 8 am, the power consumption is around 150 to 200 - 8 until 18, the
power consumption is about 250 to 300 - 18 until 23, the consumption is
around 300 to 350.

However there are some drops at the beginning of 18th Feb.~so I decided
to take the last two days as testing set so that we don't see a big
error due to these drops.

\begin{Shaded}
\begin{Highlighting}[]
\FunctionTok{filter}\NormalTok{(df,power}\SpecialCharTok{==}\DecValTok{0}\NormalTok{)}
\end{Highlighting}
\end{Shaded}

\begin{verbatim}
## # A tibble: 11 x 3
##    time                power  temp
##    <dttm>              <dbl> <dbl>
##  1 2010-02-18 00:00:00     0  12.8
##  2 2010-02-18 00:15:00     0  12.2
##  3 2010-02-18 00:30:00     0  12.2
##  4 2010-02-18 00:45:00     0  12.2
##  5 2010-02-18 01:00:00     0  12.2
##  6 2010-02-18 01:15:00     0  11.7
##  7 2010-02-18 01:30:00     0  11.7
##  8 2010-02-18 01:45:00     0  11.7
##  9 2010-02-18 02:00:00     0  11.7
## 10 2010-02-18 02:15:00     0  11.1
## 11 2010-02-18 02:30:00     0  11.1
\end{verbatim}

\#Spliting dataset

\begin{Shaded}
\begin{Highlighting}[]
\FunctionTok{library}\NormalTok{(ggplot2)}
\CommentTok{\# split the last two days for testing models }
\NormalTok{df\_train }\OtherTok{=}\NormalTok{ df[df}\SpecialCharTok{$}\NormalTok{time }\SpecialCharTok{\textless{}=} \FunctionTok{as.POSIXct}\NormalTok{(}\StringTok{"2/16/2010 23:45"}\NormalTok{, }\AttributeTok{format =} \StringTok{"\%m/\%d/\%Y \%H:\%M"}\NormalTok{),]}
\NormalTok{df\_test }\OtherTok{=}\NormalTok{ df[df}\SpecialCharTok{$}\NormalTok{time }\SpecialCharTok{\textgreater{}} \FunctionTok{as.POSIXct}\NormalTok{(}\StringTok{"2/16/2010 23:45"}\NormalTok{, }\AttributeTok{format =} \StringTok{"\%m/\%d/\%Y \%H:\%M"}\NormalTok{),]}
\NormalTok{x\_train }\OtherTok{=} \FunctionTok{ts}\NormalTok{(df\_train}\SpecialCharTok{$}\NormalTok{power, }\AttributeTok{start =} \FunctionTok{c}\NormalTok{(}\DecValTok{1}\NormalTok{,}\DecValTok{5}\NormalTok{),}\AttributeTok{frequency =} \DecValTok{24}\SpecialCharTok{*}\DecValTok{4}\NormalTok{)}
\end{Highlighting}
\end{Shaded}


\end{document}
